% BHU PYQ -- Transcribed & Corrected
% Paper: BOTMJ21 / BOTMN21: General Microbiology
% Exam: B.Sc. (Hons) Semester III (NEP) Examination 2024-25
% Category: Major / Minor Course
% Department: Botany

\documentclass[11pt,a4paper]{article}
\usepackage[a4paper,top=2.5cm,bottom=2.5cm,left=2.8cm,right=2.8cm,headheight=14pt]{geometry}
\usepackage{amsmath,amssymb}
\usepackage[T1]{fontenc}
\usepackage[utf8]{inputenc}
\usepackage{lmodern,microtype,fancyhdr,enumitem,hyperref,lastpage,parskip}

\hypersetup{
    colorlinks=true,
    urlcolor=black,
    linkcolor=black,
    pdftitle={BOTMJ21 / BOTMN21: General Microbiology -- BHU Sem III 2024-25 (NEP)}
}

\pagestyle{fancy}
\fancyhf{}
\renewcommand{\headrulewidth}{0.4pt}
\renewcommand{\footrulewidth}{0.4pt}
\fancyhead[L]{\small   \textit{Banaras Hindu University}}
\fancyhead[R]{\small   \textit{BOTMJ21 / BOTMN21: General Microbiology}}
\fancyfoot[C]{\small Page  \thepage\ of \pageref{LastPage}}


\newlist{parts}{enumerate}{2}
\setlist[parts,1]{label=(\alph*),leftmargin=*,itemsep=4pt,topsep=2pt}

\newcommand{\pts}[1]{\hfill   \textit{[#1]}}


\begin{document}


\begin{center}
  {\large\bfseries BANARAS HINDU UNIVERSITY}\\[2pt]
  {
\normalsize B.Sc. (Hons) Semester III (NEP) Examination 2024-25}\\[6pt]
  \rule{0.8\linewidth}{0.5pt}\\[6pt]
  {\large\bfseries BOTANY}\\[4pt]
  {\large\bfseries Paper No. BOTMJ21 / BOTMN21 : General Microbiology}\\[4pt]
  {
\normalsize Time: 2:30 Hours\hfill Maximum Marks: 70}
\end{center}

\rule{\linewidth}{0.4pt}

\smallskip



\noindent (Write your Roll No. at the top immediately on the receipt of this question paper)

\smallskip



\noindent   \textit{Note: Attempt five questions including Question 1 which is compulsory. The figures in the right-hand margin indicate marks.}

\bigskip


\begin{enumerate}[label=   \textbf{\arabic*.}]

    \item Answer all the following in not more than 50 words each:\pts{7 $ \times$ 2 = 14}
    
\begin{parts}
        \item Who developed postulates to prove that a microorganism causes a specific disease?
        \item What did Joseph Lister contribute to microbiology?
        \item What are biodegradation and biotransformation?
        \item What is active transport?
        \item What is the bacterial capsule?
        \item What is the biological source of the antibiotic tetracycline?
        \item What are Halophiles?
    \end{parts}

    \bigskip

    \item A bacterial colony isolated and subsequently purified from the arid desert ecosystem of Rajasthan exhibited significant morphological heterogeneity upon detailed observation. What comprehensive set of parameters can be employed to accurately delineate its taxonomic identity and how can these approaches help in resolving the presence of multiple morphotypes within a single isolate?\pts{14}

    \bigskip

    \item Write short notes on any two of the following:\pts{2 $ \times$ 7 = 14}
    
\begin{parts}
        \item Write about phenotypic characteristics used to identify microorganisms.
        \item In which scenario is DNA fingerprinting an effective tool for identifying microorganisms, and under what circumstances might it fail or be inadequate as a molecular marker?
        \item What are xenobiotics? Discuss the role of microbes in their degradation with suitable examples.
    \end{parts}

    \bigskip

    \item What is the difference between T4 phage and Lambda ($\lambda$) bacteriophage? Describe the life cycles of these bacteriophages with appropriate diagrams.\pts{14}

    \bigskip

    \item Answer any two of the following:\pts{2 $ \times$ 7 = 14}
    
\begin{parts}
        \item Differentiate between the cell walls of Gram-positive and Gram-negative bacteria.
        \item Write about Endospores.
        \item Write about Thermophiles.
    \end{parts}

    \bigskip

    \item How are antibiotics classified? Discuss the mechanisms involved in antibiotic resistance.\pts{14}

    \bigskip

    \item Write short notes on any two of the following:\pts{2 $ \times$ 7 = 14}
    
\begin{parts}
        \item Nucleoid
        \item Bacterial flagella
        \item Psychrophiles
    \end{parts}

\end{enumerate}

\end{document}
